\chapter[Equazione del trasporto]{Tecniche e metodi per la soluzione dell'equazione per il trasporto dei neutroni 
\label{cap2}}
%%%%%%%%%%%%%%%%%%%%%%%%%%%%%%%%%%%%%%%%%%%%%%%%%%%%%%%%%%%%%%%%%%%%%%%%%%%
In questo capitolo saranno affrontate le principali equazioni della neutronica 
che governano la fisica del reattore. Le caratteristiche della fisica del 
reattore associate all'interazione dei neutroni con tutti i componenti materiali 
di un reattore nucleare, influenza e in alcuni casi determina la progettazione e 
le operazioni del sistema nucleare considerato.
Uno dei più importanti parametri è il tasso di reazione con isotopi fissili che 
possono essere usati nella valutazione del rateo di fissione e la conseguente 
produzione di calore all'interno del nocciolo del reattore. Questa sorgente di 
calore determina, a sua volta, la potenza uscente dal reattore e influenza 
fortemente la temperatura dei materiali che spesso determinano le condizioni di 
sicurezza del reattore.
% % % Il tasso di reazione è determinato da una specifica sezione d'urto la quale è la 
% % % probabilità che diverse reazioni abbiano luogo quando i neutroni interagiscono 
% % % con i nuclei del materiale investito dal flusso neutronico.
Nel corso di questo capitolo saranno affrontati in un primo momento i 
concetti generali della teoria del trasporto, per assicurare una giusta 
comprensione dei concetti che portano alla definizione della equazione di 
Boltzmann per il trasporto dei neutroni.
Successivamente saranno introdotte le principali approssimazioni fisiche, usate 
nei modelli computazionali, per la risoluzione dell'equazione del trasporto 
semplificata.  
I codici computazionali che risolvono il flusso neutronico stazionario si 
possono dividere in deterministici e stocastici. 
I metodi stocastici (MonteCarlo) risolvono il flusso neutronico simulando il 
trasporto di particelle che risolve l'equazione del 
trasporto di Boltzmann mediante diversi flussi collisionali. 
Le simulazioni di particelle MonteCarlo usano un 
algoritmo, il quale assume decisioni riguardo la vita delle particelle che sono 
determinate secondo una densità di probabilità matematica.
I metodi deterministici invece coinvolgono la suddivisione numerica di variabili 
indipendenti di spazio, energia e direzione in sottodomini spaziali
con la  conseguente riformulazione dell'equazione di Boltzmann in un insieme di 
equazioni discretizzate basate sugli elementi finiti o volumi finiti. 
La risoluzione dell'insieme di equazioni discretizzate richiede tecniche e 
metodi computazionali iterativi per ottenere la convergenza della soluzione 
per ottenere il flusso neutronico, da cui deriva la determinazione di altri 
importanti parametri come la costante di criticit\`a  e la potenza neutronica.

Saranno affrontati il metodo delle caratteristiche basato sul 
concetto di traiettoria del neutrone, metodi di discretizzazione angolare quale 
il metodo $P_N$ secondo lo sviluppo in armoniche sferiche e il metodo $S_N$ 
secondo la scomposizione in formule di quadratura ed infine il metodo delle 
probabilità di collisione che deriva dall'approssimazione multigruppo 
dell'equazione integrale di Boltzmann.
In particolare i libri di riferimento da cui sono state elaborate le informazioni 
sono contenute nelle referenze \cite{lamarsh}, \cite{stacey} e soprattutto 
\cite{hebertreact}. 

\section{Teoria del trasporto}
In questa sezione saranno enunciati i concetti più importanti che descrivono la 
teoria del trasporto. Lo sviluppo della teoria del trasporto deriva dalla 
necessità di conoscere nei dettagli la distribuzione in spazio, energia e tempo 
dei neutroni all'interno dei reattori nucleari. Il trasporto dei neutroni in 
mezzi moderanti o moltiplicanti costituisce una parte della meccanica 
statistica di non equilibrio termodinamico, che studia in generale l'evoluzione 
dinamica dei sistemi formati da molte particelle. Il comportamento dinamico di 
un sistema formato da un elevato numero di particelle, quali possono essere i 
neutroni all'interno di un reattore nucleare, viene studiato attraverso la 
definizione dei processi di trasporto, processi irreversibili che 
caratterizzano la teoria del trasporto \cite{molin}.
Pertanto in fisica del reattore è di estremo interesse l'equazione del trasporto 
di Boltzmann, da cui derivano altre importanti equazioni, come l'equazione della 
diffusione opportunamente ricavata attraverso alcune ipotesi con lo scopo di 
semplificare il problema.


\subsection{Spazio delle fasi e densità di probabilità}
Un sistema dinamico è definito come un insieme costituito da $N$ particelle 
interagenti tra loro, quali possono essere neutroni, atomi o molecole, che per 
semplicità sono considerate tutte uguali e ciascuna di esse ha $l$ gradi 
di libertà. 
Sono descritte tre tipi di variabili che servono per caratterizzare il 
comportamento del sistema; quantità riguardante le particelle stesse (come per 
esempio la massa), parametri esterni che determinano le condizioni al 
contorno e le forze esterne, ed infine le cosiddette variabili dinamiche che 
descrivono il moto delle particelle.

Un particolare insieme di queste ultime variabili all'istante di tempo $t$ 
costituisce lo \textit{stato dinamico} o \textit{fase} del sistema al tempo $t$.
Infatti lo stato del sistema è completamente determinato dalla conoscenza della 
posizione $\vec{r}$ e della velocità $\vec{v}$ di ciascuna particella, e 
pertanto è necessario conoscere $6N$ numeri. 
In ogni caso questi $6N$ variabili possono essere rappresentati da un punto in 
uno spazio a $6N$ dimensioni, denominato spazio delle fasi; tuttavia la 
conoscenza esatta di tale numero è impossibile sia per l'elevato numero di 
particelle considerate sia per indeterminazioni quantistiche.
Per questo motivo si ricorre a metodi e concetti di natura statistica; infatti 
si considerano un elevato numero di sistemi uguali (ovvero stesse particelle e 
stesso numero) che differiscono tra loro solo per il micro-stato che hanno ad un 
certo istante. In questo modo questi insiemi possono essere rappresentati da una 
nube di punti nello spazio delle fasi. 
Si introduce la funzione densità di probabilità come
\begin{equation}
D_N (\vec{r_1},..,\vec{r_N}, \vec{v_1},..,\vec{v_N},t) \,,
\end{equation}
che rappresenta la probabilità di trovare lo stato del sistema in un determinato 
punto dello spazio delle fasi $\Gamma$. Infatti tale funzione $D_N$ contiene il 
massimo delle informazioni che si possono avere sul sistema.

Si introduce lo spazio delle fasi $\mu$, lo spazio a 6 dimensioni in cui lo stato 
dinamico di una particella è determinato dalle sue tre coordinate spaziali e 
dalle sue tre componenti di velocità. In questo modo, considerando tutte le 
particelle, che compongono il sistema, indistinguibili tra di loro è possibile 
introdurre la funzione di distribuzione per una particella
\begin{equation}
f_1=f(\vec{r_1},\vec{v_1},t)=\int N D_N(\vec{r_1},..,\vec{r_N}, 
\vec{v_1},..,\vec{v_N},t)\: d\vec{r_2}..d\vec{r_N}d\vec{v_2}..d\vec{v_N} \,,
\end{equation}
chiamata anche \textit{funzione di distribuzione semplice} poiché la posizione e la 
velocità di ogni data particella non è correlata a quella di un altra. Il 
principio di non correlazione equivale ad ipotizzare che funzione di densità di 
probabilità di ogni particella all'interno del sistema non è influenzata dalla 
funzione di densità di probabilità di quella di un'altra particella.

E' importante anche introdurre la funzione di distribuzione doppia, in cui due 
particelle sono correlate tra di loro,  semplicemente si ha
\begin{equation}
f_2 =\frac{N-1}{N}f_1(1)f_1(2) \,,
\end{equation}
dove $f_1(1)$ rappresenta la funzione di distribuzione semplice corrispondente 
alla particella 1, mentre $f_1(2)$ alla particella 2.


\subsection{Equazione del trasporto}
L'equazione di Louville parte dall'ipotesi che il volumetto nello spazio delle 
fasi $\Gamma$ si comporti come un fluido incomprimibile sotto l'azione di una 
hamiltoniana, ossia
\begin{equation}
\label{louville}
\frac{\partial D_N}{\partial t} + \sum_{i=1}^N \vec{v_i} \cdot \frac{\partial 
D_N}{\partial \vec{r_i}} + \sum_{i=1}^N \frac{\vec{F_i}+\sum_{i=1}^N 
\vec{F_{ij}}}{m} \cdot \frac{\partial D_N}{\partial \vec{v_i}} = 0 \,.
\end{equation} 
Da questa equazione scritta per la funzione di densità di probabilità, che 
governa il moto di una nube di sistemi dinamici, si può integrare ottenendo 
un sistema di equazioni, denominato gerarchia di equazioni di "BBGKY" per la 
funzione di distribuzione semplice. L'idea alla base di questo sistema di 
equazioni è che le quantità di interesse pratico non dipendono dalla densità di 
probabilità nello spazio delle fasi $\Gamma$ ma solo dalla probabilità di una 
singola particella di essere in un certo elemento dello spazio delle fasi $\mu$.
Quindi integrando dall'equazione \eqref{louville} e sostituendo il termine di 
collisione che si ottiene, si può scrivere l'equazione di Boltzmann con 
collisione
\begin{equation}
\label{boltzmann}
\frac{\partial f_1}{\partial t} + \vec{v_i} \cdot \frac{\partial f_1}{\partial 
\vec{r_1}} + \frac{\vec{F_1}}{m} \cdot \frac{\partial f_1}{\partial \vec{v_1}} = 
\biggl(\frac{\partial f_1}{\partial t}\biggr)_{coll} = \bigl[\Gamma (+)-\Gamma 
(-)\bigr] \,,
\end{equation}
in cui è stata introdotta nel termine di collisione l'ipotesi di urti binari 
(ovvero di durata praticamente nulla), secondo un termine di guadagno e uno di 
perdita.

\section{Equazione di Boltzmann per il trasporto di neutroni}
L'equazione del trasporto è introdotta per descrivere una popolazione di 
particelle neutre, come i neutroni, in un certo dominio chiuso, sotto condizioni 
transitorie o stazionarie.
Come già detto, l'equazione del trasporto descrive il comportamento statistico 
di un popolazione di particelle; in particolare l'esatto numero di particelle 
per unità di volume è in continuo mutamento con il tempo, e anche sotto 
condizioni stazionarie, il numero di densità di particelle oscilla attorno un 
valore medio dell'equazione stazionaria del trasporto.
Saranno enunciati diversi approcci per risolvere l'equazione del trasporto, 
ovvero il metodo delle armoniche sferiche, il metodo della probabilità di 
collisione, il metodo delle ordinate discrete e il metodo delle caratteristiche.
Partendo dalle equazioni generali del trasporto, la distribuzione dei neutroni 
in spazio $\vec{r}$, in direzione angolare $\vec{\Omega}$ e in energia $E$ 
all'istante di tempo $t$ viene descritta secondo una particolare funzione di 
distribuzione $n(\vec{r},\vec{\Omega},E,t)$, denominata funzione di densità 
neutronica o più genericamente densità di popolazione. 
Il numero atteso di neutroni, considerati in un elemento di volume $d\vec{r}$ 
alla posizione $\vec{r}$, che si muovono in un cono di direzione $d\vec{\Omega}$ 
secondo la direzione $\vec{\Omega}$ e che abbiano energia compresa tra $E$ e 
$E+dE$ in un certo instante di tempo, è rappresentato da
\begin{equation}
n(\vec{r},\vec{\Omega},E,t)\, d\vec{r}d\vec{\Omega}dE \,.
\end{equation}
Integrando rispetto l'angolo solido $d\vec{\Omega}$, si ottiene il numero atteso 
di neutroni medio nel volumetto $d\vec{r}$ in un intorno di $\vec{r}$ 
all'istante $t$ con energia in $dE$ in un intorno di $E$ come
\begin{equation}
N(\vec{r},E,t) \: d\vec{\Omega} = 
d\vec{r}dE\int_{4\pi}n(\vec{r},\vec{\Omega},E,t) \: d\vec{\Omega} \,.
\end{equation}
La variabile maggiormente usata in fisica del reattore risulta essere tuttavia 
il flusso scalare $\phi$ e il flusso angolare $\psi$ definiti come
\begin{equation}
\phi(\vec{r},E,t) = v N(\vec{r},E,t) \,,
\end{equation}
e
\begin{equation}
\psi(\vec{r},\vec{\Omega},E,t)= v n(\vec{r},\vec{\Omega},E,t) \,,
\end{equation}
dove $v=|\vec{v}|$ rappresenta il modulo della velocità che possiede il 
neutrone. Inoltre è giusto sottolineare che $\vec{\Omega}$ rappresenta il 
versore dell'angolo solido della velocità $\vec{v}$ definito come
\begin{equation}
\vec{\Omega}=\frac{\vec{v}}{|\vec{v}|} \,.
\end{equation}

Il flusso angolare definisce il massimo delle informazione riguardo la 
popolazione neutronica. Tuttavia questa quantità per come è definita risulta 
difficilmente maneggiabile durante i calcoli e sarà necessario apportare delle 
semplificazioni fisiche e matematiche al fine di ricavare la soluzione con meno 
difficoltà.
Attraverso la definizione di flusso è possibile definire il rateo di interazioni 
$R$, già definita in \eqref{tasso_di_reazione}, relativa alla sezione d'urto 
macroscopica $\Sigma$
\begin{equation}
R=v\Sigma(\vec{r}.E)N(\vec{r},E,t)=\Sigma(\vec{r}.E)\phi(\vec{r},E,t) \,.
\end{equation}
 
Un'altra importante distribuzione, correlata alla definizione di flusso 
angolare, è il vettore di corrente angolare $\vec{J}(\vec{r},\vec{\Omega},E,t)$ 
definito come
\begin{equation}
\vec{J}(\vec{r},\vec{\Omega},E,t)=v \vec{\Omega} 
n(\vec{r},\vec{\Omega},E,t)=\vec{\Omega} \psi(\vec{r},\vec{\Omega},E,t) \,.
\end{equation}
Si nota che la quantità
\begin{equation}
\vec{J}(\vec{r},\vec{\Omega},E,t)d\vec{A}dEd\vec{\Omega} \,,
\end{equation}
rappresenta il numero atteso di neutroni che transitano attraverso l'area 
$d\vec{A}$ in un intorno di $\vec{r}$ all'istante $t$, con energia $E$ in $dE$ e 
velocità nell'angolo solido $d\vec{\Omega}$ nel tempo $t$.
In analogia al flusso integrato, si può definire la corrente integrata sulla 
distribuzione angolare
\begin{equation}
\vec{J}(\vec{r},E,t)=\int_{4\pi}\vec{J}(\vec{r},\vec{\Omega},E,t)\, 
d\vec{\Omega}=\vec{\Omega}\,\phi(\vec{r},E,t) \,.
\end{equation}

Secondo l'ipotesi di assenza di interazioni tra particella e particella, ovvero 
tra neutrone e neutrone, è possibile enunciare la forma lineare dell'equazione 
di Boltzmann. In ambito reattoristico tale ipotesi è ben supportata dall'idea 
che l'interazione tra neutrone e nucleo sia molto più probabile dell'interazione 
tra neutrone e neutrone, in quanto le dimensioni del neutrone sono più piccoli 
di qualche ordine di grandezza rispetto al nucleo.

L'equazione del trasporto è ottenuta come una relazione di bilancio tra 
particelle in un volumetto di controllo $V$ dello spazio delle fasi. In questo 
caso invece dell'energia viene utilizzata l'equivalente modulo di velocità, 
semplicemente secondo le leggi della fisica di $v=\sqrt{2E/m}$ dove $m$ è la 
massa della particella considerata, in questo caso il neutrone che ha una massa di 1.67$\cdot$10$^{-27}$kg. Questo permette uno sviluppo più logico dal 
punto di visto concettuale ma formalmente risulta equivalente.
Il bilancio, tra termini di guadagno e di perdita, risulta così strutturato:
\begin{itemize}
\item[-] Variazione di neutroni in $V$ durante il $\Delta t$
\begin{equation}
\int_V [n(\vec{r},\vec{\Omega},v,t+\Delta t)-n(\vec{r},\vec{\Omega},v,t)]\, 
d\vec{r}dvd\vec{\Omega} \,,
\end{equation}
\item[-] Numero netto di particelle uscenti da $V$ durante il $\Delta t$ 
ottenuto integrando la corrente di particelle sopra la superficie $\partial V$, 
dove $\hat{n}$ è il versore normale uscente a $\partial V$ in cui 
successivamente è stato introdotto il teorema della divergenza
\begin{equation}
\int_{\partial V} (\vec{\Omega} \cdot 
\hat{n})\psi(\vec{r},\vec{\Omega},v,t+\Delta t)\, d\vec{r}dvd\vec{\Omega}\Delta 
t=\int_V \nabla \cdot \Omega \psi(\vec{r},\vec{\Omega},v,t+\Delta t)\, 
d\vec{r}dvd\vec{\Omega}\Delta t \,,
\end{equation} 
\item[-] Numero di interazioni in $d\vec{r}$ durante il $\Delta t$, dove 
$\Sigma$ è la sezione d'urto macroscopica totale considerata indipendente da 
$\vec{\Omega}$ e $t$
\begin{equation}
\int_V \Sigma(\vec{r},v)vN(\vec{r},\vec{\Omega},v,t)]\, 
d\vec{r}dvd\vec{\Omega}\Delta t \,,
\end{equation}
\item[-] Numero di particelle generate in $d\vec{r}$ durante il $\Delta t$, dove 
$Q(\vec{r},\vec{\Omega},v,t)$ rappresenta la sorgente
\begin{equation}
\int_V Q(\vec{r},\vec{\Omega},v,t)\, d\vec{r}dvd\vec{\Omega}\Delta t \,.
\end{equation}

\end{itemize}
Assemblando questi quattro termini si ottiene un equazione del bilancio, dove la 
variazione di particelle all'interno del volumetto di controllo $V$ è 
rappresenta tra il bilancio tra i neutroni quelli uscenti e entranti,
\begin{multline}
\frac{n(\vec{r},\vec{\Omega},v,t+\Delta t) - n(\vec{r},\vec{\Omega},v,t)}{\Delta 
t}= \\= - \nabla \cdot \Omega \psi(\vec{r},\vec{\Omega},v,t) 
-\Sigma(\vec{r},v)[vn(\vec{r},\vec{\Omega},v,t)]+ \\
+Q(\vec{r},\vec{\Omega},v,t) \,.
\end{multline}
Ponendo il limite di $\Delta t \rightarrow 0$, e introducendo il flusso angolare 
come variabile indipendente si ottiene l'equazione del trasporto nella sua forma 
differenziale
\begin{equation}
\frac{1}{v}\frac{\partial \psi}{\partial t}(\vec{r},\vec{\Omega},v,t)+\Omega 
\cdot \nabla \psi(\vec{r},\vec{\Omega},v,t)+ 
\Sigma(\vec{r},v)\psi(\vec{r},\vec{\Omega},v,t)=Q(\vec{r},\vec{\Omega},v,t) \,.
\end{equation}
Nelle condizioni stazionarie l'equazione si semplifica come segue
\begin{equation}
\vec{\Omega} \cdot \nabla \psi(\vec{r},\vec{\Omega},v)+ 
\Sigma(\vec{r},v)\psi(\vec{r},\vec{\Omega},v)=Q(\vec{r},\vec{\Omega},v) \,.
\end{equation}
La forma differenziale dell'equazione stazionaria del trasporto può essere 
ottenuta sostituendo la velocità del neutrone con l'energia come variabile indipendente
\begin{equation}
\label{eqtrasta}
\vec{\Omega} \cdot \nabla \psi(\vec{r},\vec{\Omega},E)+ 
\Sigma(\vec{r},E)\psi(\vec{r},\vec{\Omega},E)=Q(\vec{r},\vec{\Omega},E) \,.
\end{equation}

\subsection{Densità di sorgente stazionaria}
Assumendo che le reazioni di fissione siano isotrope nel sistema di riferimento 
del laboratorio, la densità di sorgente stazionaria $Q(\vec{r},\vec{\Omega},E)$ 
può essere scritta come
\begin{equation}
\label{denssorg}
Q(\vec{r},\vec{\Omega},E)=\frac{Q^{fiss}(\vec{r},E)}{4\pi}+\int_{4\pi}d\vec{
\Omega}'\int_0^\infty \Sigma^S(\vec{r},E'\rightarrow E,\vec{\Omega}'\rightarrow 
\vec{\Omega})\phi(\vec{r},E',\vec{\Omega}')\, dE' \,,
\end{equation}
dove $\Sigma^S(\vec{r},E'\rightarrow E,\vec{\Omega}'\rightarrow \vec{\Omega})$ è 
la sezione d'urto macroscopica differenziale di scattering mentre $ 
Q^{fiss}(\vec{r},E)$ è la sorgente isotropa di fissione.

In un mezzo isotropo, la sezione d'urto di scattering è solamente in funzione 
dell'angolo di scattering
\begin{equation}
\label{sorgfiss}
\Sigma^S(\vec{r},E'\rightarrow E,\vec{\Omega}'\rightarrow 
\vec{\Omega})=\frac{1}{2\pi}\Sigma^S(\vec{r},E'\rightarrow E,\vec{\Omega}'\cdot 
\vec{\Omega}) \,.
\end{equation}
Se si assume che la sorgente isotropa di fissione sia indipendente dall'energia 
del neutrone incidente, per ogni nuclide fissionabile $i$, l'energia dei 
neutroni emessi è distribuita in accordo alla densità di probabilità conosciuta 
come spettro di fissione $\chi_i(E)$. 
Pertanto la sorgente di fissione può essere scritta come
\begin{equation}
Q^{fiss}(\vec{r},E)=\sum_{j=1}^{J^{fiss}} \chi_j(E) \int_0^\infty dE'\, \nu(E') 
\Sigma_{f,j}(\vec{r},E')\phi(\vec{r},E')   \,,
\end{equation}
dove $J^{fiss}$ rappresenta il numero totale di isotopi fissionabili, $\nu 
\Sigma_{f,j}(\vec{r},E)$ il numero di neutroni emessi per fissione.



\subsection{Equazione di Boltzmann integro-differenziale}
L'equazione del trasporto per neutroni di Boltzmann in forma 
integro-differenziale lineare descrive l'evoluzione temporale del flusso 
angolare e viene generalmente presentata nella seguente forma, attraverso la 
sostituzione della densità di sorgente in \eqref{denssorg} e in \eqref{sorgfiss} 
nell'equazione \eqref{eqtrasta},
\begin{multline}
\biggl[ \frac{1}{v} \frac{\partial}{\partial{t}} + \vec{\Omega} \cdot \nabla + 
\Sigma^T(\vec{r},E)
\biggr] \psi(\vec{r},\vec{\Omega},E,t) \\
= \int_{0}^{\infty} dE' \int_{4\pi} 
d\vec{\Omega}'\,\Sigma^S(\vec{r},{E'\rightarrow E},
{\vec{\Omega}'\rightarrow\vec{\Omega}}) 
\psi(\vec{r},\vec{\Omega}',E',t) \\
+ \frac{\chi(\vec{r},E)}{4\pi} \int_{0}^{\infty} dE'\,
\nu(E')\Sigma_f(\vec{r},E') \int_{4\pi} d\vec{\Omega}'\,
\psi(\vec{r},\vec{\Omega}',E',t) \,.
\end{multline}
L'incognita in sette variabili è il flusso angolare 
$\psi(\vec{r},\vec{\Omega},E,t)$, e per la soluzione è necessario che vengano 
definite appropriate condizioni iniziali, alle interfacce e al contorno. 
Per quanto riguarda la condizione iniziale è necessario fissare il flusso 
angolare al $t=0$ per tutte le posizioni, energie e direzioni come
\begin{equation}
\psi(\vec{r},\vec{\Omega},E,0)=\psi_0(\vec{r},\vec{\Omega},E) \,.
\end{equation} 
Le condizioni al contorno dipendono dallo specifico problema considerato, e 
quindi bisogna appunto conoscere la distribuzione del flusso angolare al 
contorno $\partial V$ che entra nel dominio spaziale
\begin{equation}
\psi(\vec{r_0},\vec{\Omega},E,t)=\beta \psi(\vec{r_0},\vec{\Omega}',E,t) \qquad 
\text{per} \quad \vec{r_0}\in \partial V,\quad\vec{\Omega}\cdot\hat{n}<0 \,.
\end{equation}
La condizione di vuoto è determinata da $\beta = 0$, ovvero dove il neutrone che 
attraversa la superficie è perso, la condizione di completa riflessione da 
$\beta=1$ mentre un valore intermedio di $\beta$ rappresenta una riflessione 
parziale. Queste condizione così espressa è detta anche condizione al contorno 
di albedo. Inoltre all'interno di un reattore nucleare sono presenti molti 
materiali diversi, e la configurazione geometrica del core può essere davvero 
complessa. Mentre le sezioni d'urto posso essere assunte funzioni continue in 
ogni materiale, all'interfaccia tra regioni spaziali differenti possono essere 
discontinue. Di conseguenza l'equazione del trasporto neutronico deve essere 
applicata in entrambi i lati dell'interfaccia, imponendo condizioni di 
continuità sul flusso angolare.

Se si considera l'equazione del trasporto di Boltzmann stazionaria si può 
introdurre il concetto di criticità di un sistema moltiplicante che è legato 
intrinsecamente al concetto di stazionarietà
\begin{multline}
\label{eq_boltzmann_staz_estesa}
\biggl[\vec{\Omega} \cdot \nabla + \Sigma^T(\vec{r},E)
\biggr] \psi(\vec{r},\vec{\Omega},E) \\
= \int_{0}^{\infty} dE' \int_{4\pi} 
d\vec{\Omega}'\,\Sigma_s(\vec{r},{E'\rightarrow E},
{\vec{\Omega}'\rightarrow\vec{\Omega}}) 
\psi(\vec{r},\vec{\Omega}',E') \\
+ \frac{\chi(\vec{r},E)}{4\pi
k_{eff}} \int_{0}^{\infty} dE'\,
\nu(E')\Sigma_f(\vec{r},E') \int_{4\pi} d\vec{\Omega}'\,
\psi(\vec{r},\vec{\Omega}',E') \,.
\end{multline}
Si possono introdurre operatori per esprimere l'equazione come un problema agli autovalori ed
evidenziare come l'autovalore sia
oppurtanamento inserito per ottenere la criticità del sistema moltiplicativo, 
ovvero
\begin{equation}
\begin{split}
\vec{T}\psi & = \biggl[\vec{\Omega} \cdot \nabla + \Sigma^T(\vec{r},E)
\biggr] \psi(\vec{r},\vec{\Omega},E), \\
\vec{S}\psi &= \int_{0}^{\infty} dE' \int_{4\pi} 
d\vec{\Omega}'\,\Sigma_s(\vec{r},{E'\rightarrow E},
{\vec{\Omega}'\rightarrow\vec{\Omega}}) 
\psi(\vec{r},\vec{\Omega}',E'),\\
\vec{F}\psi &= \frac{\chi(\vec{r},E)}{4\pi} \int_{0}^{\infty} dE'\,
\nu(E')\Sigma_f(\vec{r},E') \int_{4\pi} d\vec{\Omega}'\,
\psi(\vec{r},\vec{\Omega}',E') \,,
\end{split}
\end{equation}
dove gli operatori sono definiti, rispettivamente, di trasporto $\vec{T}$, di 
scattering $\vec{S}$ e di fissione $\vec{F}$. Nella progettazione di reattori 
nucleari si cerca la soluzione di un problema agli autovalori ottenuto inserendo 
nell’ equazione un parametro di valutazione per lo scostamento dalla condizione 
di criticità.
Introducendo l'autovalore $\frac{1}{k_{eff}}$ in forma operatoriale il problema agli autovalori diventa
\begin{equation}
(\vec{T}-\vec{S})\psi = \frac{1}{k_{eff}} \vec{F}\psi \,.
\end{equation}
La costante moltiplicativa effettiva $k_{eff}$ rappresenta quindi la criticità 
del sistema, e se $k_{eff}<1$ il sistema sarà sottocritico ovvero la 
popolazione neutronica andrà estinguendosi con il tempo (cioè la produzione 
di neutroni è inferiore alle perdite) mentre se $k_{eff}>1$ il sistema sarà 
sovracritico e la popolazione neutronica aumenterà ad ogni passo temporale.
La condizione di criticità è invece data da $k_{eff}=1$, in cui la popolazione 
rimane stabile, la produzione di neutroni è compensata dalle perdite e 
la reazione a catena di fissione è in grado di autosostenersi senza estinguersi 
o esplodere.
Infatti nel caso reale di reattore nucleare, questo parametro viene 
continuamente aggiustato attraverso inserzioni di reattività attraverso 
l'inserimento di barre di controllo o attraverso l'iniezione di veleni 
neutronici.


\subsection{Metodo delle caratteristiche \label{metcaratt}}
Attraverso la rappresentazione dell'inversione dell'operatore del trasporto 
lungo la direzione di volo dei neutroni $\vec{\Omega}$, è possibile riformulare 
il problema del trasporto in forma integrale. Questa procedura è conosciuta come 
metodo delle caratteristiche \cite{askew}.

La forma caratteristica dell'equazione del trasporto corrisponde ad un'
integrazione dell'operatore di streaming $\vec{\Omega}\cdot\nabla\psi$ lungo la 
caratteristica, in linea diretta lungo la direzione $\vec{\Omega}$ 
corrispondente alla traiettoria della particella. 
Ad ogni istante di tempo del suo moto, la particella è assunta essere alla 
distanza $s$ dalla posizione di riferimento $\vec{r_0}$ sulla propria 
caratteristica, che è l'intersezione con l'intera superficie $\partial V$, così 
che l'attuale posizione di ogni punto lungo $\vec{\Omega}$ può essere espressa 
come
\begin{equation}
\vec{r}= (x(s),y(s),z(s))=(x_0+s\Omega_x, y_0+s\Omega_y, 
z_0+s\Omega_z)=\vec{r_0}+s\,\vec{\Omega} \,,
\end{equation}
dove $s$ è la lunghezza del segmento lungo la caratteristica. In altre parole si 
è ottenuta la parametrizzazione di $\vec{r}$ rispetto a una posizione di 
riferimento $\vec{r_0}$.

Questo segmento può essere misurato in termini di libero cammino, portando alla 
definizione di cammino ottico, come segue
\begin{equation}
\label{cammino_ottico}
\tau (s,E) =\int_0^s ds'\, \Sigma(\vec{r_0}+s'\,\vec{\Omega},E) \,.
\end{equation}
Allo stesso tempo, il tempo generico $t$ può essere riferito alla velocità del 
neutrone $v$
\begin{equation}
t=t_0+\frac{s}{v} \,,
\end{equation}
dove $t_0$ è l'istante con il quale il neutrone considerato giace sul punto di 
riferimento $\vec{r_0}$.
Esprimendo il termine $\nabla\cdot\vec{\Omega}$ usando il concetto di derivata 
totale si può scrivere
\begin{equation}
\frac{d}{ds}=\Omega_x\frac{\partial}{\partial 
x}+\Omega_y\frac{\partial}{\partial y}+\Omega_z\frac{\partial}{\partial 
z}=\frac{dt}{dl}\frac{\partial}{\partial 
t}+\frac{dx}{dl}\frac{\partial}{\partial 
x}+\frac{dy}{dl}\frac{\partial}{\partial 
y}+\frac{dz}{dl}\frac{\partial}{\partial z}+\frac{1}{v}\frac{\partial}{\partial 
t} \,,
\end{equation}
dove $ds\,\vec{\Omega}$ è stato scomposto secondo un sistema di coordinate cartesiane
\begin{equation}
ds\,\vec{\Omega}=d\vec{r}=dx \, \hat{i}+dy \, \hat{j}+dz \, \hat{k} \,.
\end{equation} 
Si deduce che 
\begin{equation}
\frac{d}{ds}=(\vec{\Omega}\cdot \hat{i})\frac{\partial}{\partial 
x}+(\vec{\Omega}\cdot \hat{j})\frac{\partial}{\partial y}+(\vec{\Omega}\cdot 
\hat{k})\frac{\partial}{\partial z}+\frac{1}{v}\frac{\partial}{\partial t} \,.
\end{equation}
Con questa definizione è possibile riformulare l'equazione del trasporto come
\begin{multline}
\frac{d}{ds}\psi\biggl(\vec{r_0}+s\,\vec{\Omega},\vec{\Omega},E,t_0+\frac{s}{v}
\biggr)\\
+\Sigma\biggl(\vec{r_0}+s\,\vec{\Omega},\vec{\Omega},E,t_0+\frac{s}{v}
\biggr)\psi\biggl(\vec{r_0}+s\,\vec{\Omega},\vec{\Omega},E,t_0+\frac{s}{v}
\biggr)\\
=Q\biggl(\vec{r_0}+s\,\vec{\Omega},\vec{\Omega},E,t_0+\frac{s}{v}\biggr) \,.
\end{multline}
Quest'ultima, che è un'equazione differenziale lineare del primo ordine, può essere 
integrata da $0$ a $\infty$ così da ottenere
\begin{equation}
\psi(\vec{r},\vec{\Omega},E,t)=\int_0^\infty ds\, e^{-\tau 
(s,E)}Q\biggl(\vec{r_0}+s\,\vec{\Omega},\vec{\Omega},E,t_0+\frac{s}{v}\biggr) 
\,.
\end{equation}
Se il dominio è finito, è possibile integrare solo lungo il valore $s$ che 
corrisponde al valore $r'$ dentro al dominio dove il flusso angolare al 
$(\vec{r_0},t_0)$ è assunto come flusso al contorno, ottenendo
\begin{equation}
\psi(\vec{r},\vec{\Omega},E,t)= \psi(\vec{r_0},\vec{\Omega},E,t_0)e^{-\tau(l,E)}
+\int_0^l ds\, e^{-\tau 
(s,E)}Q\biggl(\vec{r_0}+s\,\vec{\Omega},\vec{\Omega},E,t_0+\frac{s}{v}\biggr) 
\,.
\end{equation}

Allo stesso modo è possibile definire la trasformazione caratteristica 
dell'equazione stazionaria del trasporto nella seguente forma
\begin{equation}
\vec{\Omega}\cdot\nabla \psi 
(\vec{r_0}+s\vec{\Omega},\vec{\Omega},E)+\Sigma^T(\vec{r_0}+s\vec{\Omega},E) 
\psi(\vec{r_0}+s\vec{\Omega},\vec{\Omega},E) = 
Q(\vec{r_0}+s\vec{\Omega},\vec{\Omega},E) \,.
\end{equation}
Successivamente il problema pu\`o essere semplificato assumendo la sola dipendenza da $s$, 
relativa alla posizione di riferimento $\vec{r_0}$ e dalla direzione angolare 
$\vec{\Omega}$, ottenendo così l'equazione 
\begin{equation}
\label{eqcarat}
\frac{d}{ds}\psi(s,\vec{\Omega},E)+\Sigma^T(s,E)\psi(s,\vec{\Omega},E)=Q(s,\vec{
\Omega},E) \,.
\end{equation}
Tale equazione può essere risolta attraverso l'uso di un fattore integrativo
\begin{equation}
e^{-\int^s_0 ds' \, \Sigma^T(s',E)} \,,
\end{equation}
ottenendo così la soluzione analitica dell'equazione caratteristica del 
trasporto
\begin{equation}
\label{soluz_caratteristica}
\psi(s,\vec{\Omega},E)=\psi(\vec{r_0},\vec{\Omega},E) e^{-\int^s_0 ds' \, 
\Sigma^T(s',E)} +\int_0^s ds'' \, Q(s'',\vec{\Omega},E) e^{-\int^s_{s''} ds' \, 
\Sigma^T(s',E)} \,.
\end{equation}
Dall'equazione \eqref{soluz_caratteristica} derivano importanti approssimazioni, 
come per esempio il metodo della probabilità di collisione.



\section{Discretizzazione dell'equazione del trasporto}
La soluzione numerica dell'equazione del trasporto stazionaria richiede la 
discretizzazione delle variabili che descrivono le caratteristiche (posizione, 
direzione angolare, e energetica) dei neutroni nello spazio delle fasi. 

\subsection{Discretizzazione energetica}
La discretizzazione energetica, o approssimazione multigruppo o anche 
condensazione energetica, consiste nella partizione dell'intero intervallo 
energetico di interesse in $G$ gruppi di energia definiti come 
\begin{equation}
\Delta E_g = [E_g,E_{g-1}] \quad \text{per} \quad g=1,\dots ,G \,,
\end{equation}
dove $g=1$ corrisponde al gruppo energetico più alto, e $g=G$ a quello più 
basso.
In questo modo il flusso angolore e scalare in forma multigruppo sono esprimibili
\begin{equation}
\psi_g(\vec{r},\vec{\Omega})=\int_{E_g}^{E_{g-1}} 
dE\,\psi(\vec{r},\vec{\Omega},E) \,,
\end{equation}
\begin{equation}
\phi_g(\vec{r})=\int_{E_g}^{E_{g-1}} dE\,\phi(\vec{r},E) \,.
\end{equation}
La definizione formale della sezione d'urto mediata sui gruppi è fondata sul 
 tasso di reazione di ogni gruppo. Si pu\`o scrivere
\begin{equation}
\Sigma_g(\vec{r})=\frac{\int_{\Delta E_g} dE \, 
\Sigma(\vec{r},E)\phi(\vec{r},E)}{\phi_g(\vec{r})} \,.
\end{equation}
In questo modo le sezioni d'urto e il flusso angolare non sono più funzioni 
continue dell'energia, ma sono costanti nel dominio energetico dentro ogni gruppo 
e si rimuove dall'equazione del trasporto la dipendenza dell'energia, ottenendo 
una serie di equazioni accoppiate tra loro.
Partendo dall'equazione \eqref{eqcarat} si ottiene l'approssimazione multigruppo
\begin{equation}
\label{eq_multigrup}
\vec{\Omega}\cdot \nabla	
\psi_g(s,\vec{\Omega})+\Sigma_g^T(s)\psi_g(s,\vec{\Omega})=Q_g(s,\vec{\Omega}) 
\,,
\end{equation}
dove la sorgente è definita,
\begin{multline}
Q_g(s,\vec{\Omega})=\sum_{g'=1}^G \int_{4\pi} d\vec{\Omega}'\, \Sigma_{g' 
\rightarrow g}^S (s,\vec{\Omega}'\rightarrow 
\vec{\Omega})\psi_{g'}(s,\vec{\Omega}')\\
+\frac{\chi_g(s)}{4\pi k_{eff}}\sum_{g'=1}^G\nu \Sigma_{g'}^F (s) 
\int_{4\pi}d\vec{\Omega}'\, \psi_{g'}(s,\vec{\Omega}') \,.
\end{multline}

L'approssimazione multigruppo trasforma l'equazione del trasporto di neutroni in 
un sistema di $G$ equazioni, dove ogni gruppo di energia descrive un bilancio 
neutronico in un unico intervallo di energia.
L'accoppiamento tra questi gruppi è ottenuto attraverso l'integrale di 
scattering e il termine di fissione, in accordo con la fisica dei processi di 
collisione.

La soluzione dell'equazione del trasporto neutronico  in Eq. \eqref{soluz_caratteristica} lungo la caratteristica  in forma multigruppo risulta essere
\begin{equation}
\psi_g(s,\vec{\Omega})=\psi_g(\vec{r_0},\vec{\Omega})e^{-\int_0^s 
ds'\,\Sigma_g^T(s')}+\int_0^s ds'' \, Q_g(s'',\vec{\Omega})e^{-\int_{s''}^s 
ds'\,\Sigma_g^T(s')} \,,
\end{equation}
in cui le sono state usate le condensazioni in energia delle sezioni d'urto 
$\Sigma^T$, $\Sigma^F$ e $\Sigma^S$ e dello spettro di fissione $\chi$ definite
\begin{equation}
\Sigma_g^T(s)=\frac{\int^{E_{g-1}}_{E_g} 
dE'\,\Sigma^T(s,E')\psi(s,\vec{\Omega},E')}{\int^{E_{g-1}}_{E_g} 
dE'\,\psi(s,\vec{\Omega},E')} \,,
\end{equation}
\begin{equation}
\Sigma_g^F(s)=\frac{\int^{E_{g-1}}_{E_g} 
dE'\,\Sigma^F(s,E')\psi(s,\vec{\Omega},E')}{\int^{E_{g-1}}_{E_g} 
dE'\,\psi(s,\vec{\Omega},E')} \,,
\end{equation}
\begin{multline}
\Sigma_{g'\rightarrow g}^S(s,\vec{\Omega}'\rightarrow \vec{\Omega})= \\
=\frac{\int^{E_{g'-1}}_{E_{g'}} dE'\, \int^{E_{g-1}}_{E_g} 
dE''\,\Sigma^S(s,\vec{\Omega}'\rightarrow \vec{\Omega},E'\rightarrow 
E'')\psi(s,\vec{\Omega},E'')}{\int^{E_{g-1}}_{E_g} 
dE'\,\psi(s,\vec{\Omega}',E')} \,,
\end{multline}
\begin{equation}
\chi_{g'\rightarrow g}(s)=\frac{\int^{E_{g'-1}}_{E_{g'}} dE'\, 
\int^{E_{g-1}}_{E_g} dE''\,\chi(s,E'\rightarrow E'')\nu 
\Sigma^F(s,E')\psi(s,\vec{\Omega},E')}{\int^{E_{g-1}}_{E_g} dE'\,\nu 
\Sigma^F(s,E')\psi(s,\vec{\Omega},E')} \,.
\end{equation}
Nonostante si assume una dipendenza di $\chi$ su entrambe le energie del 
neutrone che causa fissione $g'$ e del neutrone emesso da fissione $g$, 
tipicamente viene approssimato attraverso una sommatoria su tutti i gruppi
\begin{equation}
\chi_g(s)=\sum_{g=1}^G \chi_{g'\rightarrow g}(s) \,.
\end{equation}
Ricordando, sono definite generalmente tre zone di energia: zona veloce da 10MeV a 300keV, 
zona epitermica da 300keV a 0.625eV e zona termica da 0.625eV a 0.001eV. Mentre 
nella zona veloce e epitermica lo scattering diminuisce maggiormente l'energia 
dei neutroni (down-scattering), nella zona termina il neutrone può acquisire 
energia dalle collisioni (up-scattering). 


\subsection{Discretizzazione angolare}
L'approssimazione numerica dell'equazione del trasporto richiede una 
discretizzazione angolare. I metodi più comuni per discretizzare la variabile 
angolare sono: il metodo delle armoniche sferiche ($P_N$) e il metodo delle 
ordinate discrete ($S_N$).


\subsubsection{Armoniche sferiche}

Il metodo delle armoniche sferiche, o metodo $P_N$, consiste nella 
discretizzazione angolare dell'equazione differenziale del trasporto; si basa 
infatti sull'espansione in serie di $\psi_g(\vec{r},\vec{\Omega})$ e 
$Q_g(\vec{r},\vec{\Omega})$ in armoniche sferiche, dove il pedice $g$ indica 
l'approssimazione multigruppo. Attraverso l'espansione in serie è possibile 
separare la variabili dipendenti. Le armoniche sferiche sono definite di 
seguito
\begin{equation}
f(\vec{\Omega})=\sum_{l=0}^L \frac{2l+1}{4\pi} \sum_{m=-l}^l f_l^m 
R_l^m(\vec{\Omega}) \,.
\end{equation}
dove $f(\vec{\Omega})$ è una funzione continua nella variabile angolare 
$\vec{\Omega}$ e $R_l^m(\vec{\Omega}$ la componente reale dell'armonica sferica, 
dove $f_l^m$ è definito come 
\begin{equation}
f_l^m=\int_{4\pi}d\vec{\Omega}\, R_l^m(\vec{\Omega}) f(\vec{\Omega}) \,.
\end{equation}
Tale espansione è troncata dopo alcuni termini, portando all'approssimazione 
$P_N$. Si può scrivere quindi l'espansione in armoniche sferiche
\begin{equation}
\psi_g(\vec{r},\vec{\Omega})=\sum_{l=0}^n \frac{2l+1}{4\pi}\sum_{m=-l}^l 
\psi_{g,l}^m(\vec{r})R_l^m(\vec{\Omega}) \,,
\end{equation}
\begin{equation}
Q_g(\vec{r},\vec{\Omega})=\sum_{l=0}^L\frac{2l+1}{4\pi}\sum_{m=-l}^l 
Q_{g,l}^m(\vec{r})R_l^m(\vec{\Omega}) \,,
\end{equation}
dove $R_l^m(\Omega)$ rappresenta l'armonica sferica considerata e 
$\phi_l^m(\vec{r},E)$, $Q_l^m(\vec{r},E)$ sono rispettivamente
\begin{equation}
\phi_{g,l}^m(\vec{r})= \int_{4\pi}d\vec{\Omega}\, 
R_l^m(\vec{\Omega})\psi_g(\vec{r},\vec{\Omega}) \,,
\end{equation}
\begin{equation}
Q_{g,l}^m(\vec{r})= \int_{4\pi}d\vec{\Omega}\, 
R_l^m(\vec{\Omega})Q_g(\vec{r},\vec{\Omega}) \,,
\end{equation}
dove $n$ è dispari e $L\leq n$. 
Sostituendo all'equazione del trasporto stazionaria si ottiene così un insieme 
di equazioni accoppiate tra loro, la cui risoluzione è tuttavia complessa e 
generalmente son attuate semplificazioni di geometria e fisiche per semplificare 
ulteriormente il problema matematico.

Richiamando brevemente l'equazione del trasporto multigruppo \eqref{eq_multigrup} e considerando alcune semplificazioni, si può scrivere
\begin{equation}
\vec{\Omega}\cdot\nabla \psi_g(\vec{r},\vec{\Omega})+\Sigma^T_g 
\psi_g(\vec{r},\vec{\Omega}) = 
\int_{4\pi}d\vec{\Omega}'\,\Sigma_g^S(\vec{\Omega}'\rightarrow 
\vec{\Omega})\psi_g(\vec{r},\vec{\Omega}') + Q^{fiss}_g(\vec{r},\vec{\Omega}) 
\,,
\end{equation}
dove si è scomposta la sorgente $Q_g(\vec{r},\vec{\Omega})$ in sorgente di 
scattering e sorgente di fissione.
Pertanto è stato considerato un mezzo isotropo ed omogeneo così che le sezioni 
d'urto  $\Sigma_g^T$ e $\Sigma_g^S$ dipendono solamente dal coseno dell'angolo 
di scattering, denominato anche angolo di deviazione, definito come
\begin{equation}
\vec{\Omega}'\rightarrow 
\vec{\Omega}=\vec{\Omega}'\cdot\vec{\Omega}=\cos(\theta_0)=\mu_0=\mu\mu'+\sqrt{
1-\mu^2}\sqrt{1-{\mu'}^2}cos(\phi-\phi') \,,
\end{equation}
dove $\mu=\vec{\Omega}\cdot \hat{k}$ rappresenta il coseno nel centro di massa, 
$\theta$ l'angolo polare e $\phi$ l'angolo azimutale.
Inoltre considerando una simmetria azimutale, è possibile semplificare il 
problema tridimensionale ad una sola coordinata, studiando così il problema 
unidimensionale
\begin{equation}
\mu \frac{\partial}{\partial z}\psi_g(z,\mu)+\Sigma_g^T \psi_g(z,\mu) = 
\int_{-1}^{+1}\int_0^{2\pi}d\mu'd\phi'\,\Sigma_g^S(\mu_0)\psi_g(z,\mu ') 
+Q^{fiss}_g(z,\mu) \,.
\end{equation}
In questo modo si ottiene un'equazione sviluppabile secondo i polinomi di 
Legendre, un insieme completo di funzioni ortogonali nell'intervallo $[-1,+1]$, 
che godono di proprietà di ortogonalità e relazioni di ricorrenza.
Sostituendo gli sviluppi in serie e risolvendo l'integrale si ottiene
\begin{multline}
\sum_{l=0}^\infty \frac{2l+1}{4\pi}\mu P_l(\mu) \frac{\partial}{\partial 
z}\psi_l(z)+ \sum_{l=0}^\infty \frac{2l+1}{4\pi} \Sigma^T P_l(\mu)\psi_l(z)\\
=\sum_{l=0}^\infty \frac{2l+1}{4\pi}\Sigma_l^S P_l(\mu)\psi_l(z) 
+\sum_{l=0}^\infty \frac{2l+1}{4\pi} P_l(\mu)Q^{fiss}_l(z) \,,
\end{multline}
in cui è 
stato omesso il pedice $g$ per semplicità di lettura.
Attraverso la relazione di ricorrenza
\begin{equation}
(2l+1)\mu P(\mu)=(l+1)P_{l+1}(\mu)+l P_{l-1}(\mu) \,,
\end{equation}
e semplificando i polinomi $P_l$, si ottiene un sistema di equazioni accoppiate 
tra loro come segue
\begin{equation}
\sum_{l=0}^\infty l 
\frac{d}{dz}\psi_{l-1}(z)+(l+1)\frac{d}{dz}\psi_{l+1}(z)+(2l+1)[\Sigma^T 
-\Sigma_l^S]\psi_l(z)=\sum_{l=0}^\infty Q_l^{fiss}(z) \,.
\end{equation}
Il sistema è risolto troncando l'espansione in serie ad $l=L$, che prende il 
nome di approssimazione $P_L$.
Un'ulteriore approssimazione, molto comune, è di troncare l'espansione ad $L=1$, 
ottenendo l'approssimazione $P_1$, in cui si considera lo scattering linearmente 
anisotropo (per $L=0$ lo scattering sarebbe isotropo).
Questo consente attraverso alcune sostituzioni di ricavare l'equazione della 
diffusione con correzione del trasporto.
Infatti considerando la sorgente di fissione isotropa, e sostituendo il flusso 
$\psi_0(z)=\psi(z)$ e la corrente di flusso $\psi_1(z)=J(z)$, si ricava
\begin{equation}
\label{sistema_di_diffusione}
\begin{sistema}
(\Sigma^T-\Sigma^S)\psi(z)+\displaystyle\frac{d}{dz}J(z)=Q^{fiss}(z)\\
\displaystyle\frac{d}{dz}\psi(z)+3(\Sigma^T-\bar{\mu_0}\Sigma^S)J(z)=0
\end{sistema} \,,
\end{equation}
dove $\Sigma_0^S$ e $\Sigma_1^S$ sono stati sostituiti rispettivamente da 
$\Sigma^S$ e $\bar{\mu_0}\Sigma^S$.
Sistemando le equazioni, si ottiene dalla seconda la legge di Fick nella forma
\begin{equation}
\label{corrente legendre}
J(z)= -\frac{1}{3(\Sigma^T-\bar{\mu_0}\Sigma^S)} \frac{d}{dz}\psi(z)  
-D\frac{d}{dz}\psi(z) \,,
\end{equation}
dove $D$ è definito come il coefficiente di diffusione e corrisponde, nel caso 
la sezione d'urto d'assorbimento $\Sigma^A$ sia molto più piccola della sezione 
d'urto di scattering $\Sigma^S$ alla sezione d'urto del trasporto, definita come 
$\Sigma^{tr}=\Sigma^S(1-\bar{\mu_0})$.
Quindi sostituendo \eqref{corrente legendre} in \eqref{sistema_di_diffusione} e 
ricordando che $\Sigma^A=\Sigma^T-\Sigma^S$, si ottiene 
\begin{equation}
D\frac{d^2}{dz^2}\psi(z)-\Sigma^A \psi(z) + Q^{fiss}(z)=0 \,,
\end{equation}
si ottiene l'equazione di diffusione per il caso monodimensionale. Questa 
equazione può essere risolta facilmente anche se per arrivarci son state fatte 
diverse approssimazioni rilevanti.

\subsubsection{Ordinate discrete}
Il metodo delle ordinate discrete è introdotta per approssimare l'integrale 
lungo il dominio angolare della sorgente $Q(s,\vec{\Omega})$ ed è equivalente ad 
applicare regole di quadratura per calcolare l'integrale sopra il flusso 
angolare usando una somma pesata di flussi ad uno specifico angolo dove i pesi 
$\omega_m$ sono introdotti per ogni punto di quadratura $\vec{\Omega_m}$ con $m 
\in \{1,\dots,M\}$.
In questo modo si ottiene il flusso scalare $\phi_g(s)$, considerata 
l'approssimazione multigruppo, il quale dipende solamente dalla caratteristica
\begin{equation}
\label{flusso_multigrup_ord_discrete}
\phi_g(s)=\int_{4\pi}d\vec{\Omega}'\psi_g(s,\vec{\Omega}') \approx \sum_{m=1}^M 
\omega_m \psi_g(s,\vec{\Omega_m}) \,.
\end{equation}
Allo stesso modo, si ottiene l'approssimazione per la sorgente $Q_{m,g}(s) 
\approx Q_g(s,\vec{\Omega_m})$ sostituendo il flusso angolare integrato, come 
segue
\begin{multline}
\label{sorgente_multigruppo_ord_discrete}
Q_{m,g}(s,\vec{\Omega_m})=\sum_{g'=1}^G \sum_{m'=1}^M \omega_m  \Sigma_{g' 
\rightarrow g}^S (s,\vec{\Omega_{m'}}\rightarrow 
\vec{\Omega_m})\psi_{g'}(s,\vec{\Omega_{m'}})\\
+\frac{\chi_g(s)}{4\pi k_{eff}}\sum_{g'=1}^G \sum_{m'=1}^M \omega_{m'}\nu 
\Sigma_{g'}^F (s) \int_{4\pi}d\vec{\Omega}'\, \psi_{g'}(s,\vec{\Omega_{m'}}) \,.
\end{multline}
Sostituendo questa approssimazione alla sorgente si ottiene la soluzione 
caratteristica per il flusso angolare $\psi_{m,g}(s)\approx 
\psi_g(s,\vec{\Omega_m}$ per ogni punto di quadratura $\vec{\Omega_m}$,
\begin{equation}
\psi_{m,g}(s)=\psi_{m,g}(\vec{r_0})e^{-\int_0^s ds' \, 
\Sigma_g^T(s')}+\int_0^{s_m} ds'' \, Q_{m,g}(s'')e^{-\int_{s''}^s ds' \, 
\Sigma_g^T(s')} \,.
\end{equation}
Tuttavia è possibile decomporre le equazione \eqref{flusso_multigrup_ord_discrete} e \eqref{sorgente_multigruppo_ord_discrete} secondo l'angolo azimutale 
e polare attraverso regole di quadratura.

Un'ulteriore approssimazione consiste nell'assumere che la sorgente di 
scattering, oltre a quella di fissione, sia isotropa. Tale semplificazione 
permette di esprimere la sorgente totale unicamente in termini di flusso 
scalare
\begin{equation}
\label{sorgente_isotropa_od}
Q_g(s)=\frac{1}{4\pi}\biggl(\sum_{g'=1}^G \Sigma_{g'\rightarrow 
g}^S(s)\phi_{g'}(s)+\frac{\chi_g(s)}{k_{eff}}\sum_{g'=1}^G\nu 
\Sigma_{g'}^F(s)\phi_{g'}(s)\biggr) \,.
\end{equation}

Un'altra approssimazione è assumere che la sorgente $Q_g$ sia costante 
attraverso le regioni FSRs, denominate regioni di sorgente piana (Flat Source 
Region); in questo modo si ipotizza che la sorgente non vari lungo la 
caratteristica $k$ entrando in $s'$ e uscendo in $s''$ della regione $FRS_i$
\begin{equation}
Q_{i,g}=Q_g(s')=Q_g(s'')=Q_g(s), \quad con \quad s \in [s',s''] \,.
\end{equation} 
Una descrizione più accurata della sorgente in regioni spaziali è di assumere 
una variazione lineare della sorgente $Q_g(s)$; la sorgente così varia lungo 
ogni linea caratteristica
\begin{equation}
Q_g(s)=q_{t,g,0}+q_{t,g,1}\biggl(s-\frac{l_t}{2}\biggr) \,,
\end{equation}
dove $l_t$ è la lunghezza del segmento considerato, mentre $q_{t,g,0}$ e 
$q_{t,g,1}$ sono coefficienti dipendenti dalla traiettoria che descrivono la 
sorgente.

In aggiunta alla considerazione di sorgente piana all'interno di ogni cella 
spaziale, si assume anche che le proprietà dei materiali siano costanti in ogni 
FRSs. Si possono definire le sezioni d'urto mediate sull'area del $FSR_i$ con $i 
\in \{1,\dots,I\}$ con area $A_i$, come segue
\begin{equation}
\Sigma_{i,g}^T= \frac{\int_{\vec{r}\in A_i} d\vec{r}\, 
\Sigma_g^T(\vec{r})\phi_g(\vec{r})}{\int_{\vec{r}\in A_i}d\vec{r}\, 
\phi_g(\vec{r})} \,,
\end{equation}
\begin{equation}
\Sigma_{i,g}^F= \frac{\int_{\vec{r}\in A_i} d\vec{r}\, 
\Sigma_g^F(\vec{r})\phi_g(\vec{r})}{\int_{\vec{r}\in A_i}d\vec{r}\, 
\phi_g(\vec{r})} \,,
\end{equation}
\begin{equation}
\Sigma_{i,g'\rightarrow g}^S= \frac{\int_{\vec{r}\in A_i} d\vec{r}\, 
\Sigma_{g'\rightarrow g}^S(\vec{r})\phi_{g'}(\vec{r})}{\int_{\vec{r}\in A_i} 
d\vec{r}\, \phi_{g'}(\vec{r})} \,,
\end{equation}
\begin{equation}
\chi_{i,g}= \frac{\int_{\vec{r}\in A_i} d\vec{r}\, 
\chi_g(\vec{r})}{\int_{\vec{r}\in A_i}d\vec{r}} \,.
\end{equation}
Con le approssimazioni fatte, il termine di sorgente piana $Q_{i,g}$ per la 
$FRS_i$ corrispondente con area $A_i$ è definito dal flusso scalare mediato 
sull'area $\phi_{i,g}$. Quindi dalla \eqref{sorgente_isotropa_od} si ottiene
\begin{equation}
Q_{i,g}(s)=\frac{1}{4\pi}\biggl(\sum_{g'=1}^G \Sigma_{i,g'\rightarrow 
g}^S(s)\phi_{i,g'}(s)+\frac{\chi_g(s)}{k_{eff}}\sum_{g'=1}^G\nu 
\Sigma_{i,g'}^F(s)\phi_{i,g'}(s)\biggr) \,,
\end{equation}
dove $\phi_{i,g}$ è definito come
\begin{equation}
\label{flusso_scalare_FSR}
\phi_{i,g}=\frac{\int_{\vec{r}\in A_i} d\vec{r}\, 
\phi_g(\vec{r})}{\int_{\vec{r}\in A_i} d\vec{r}} \,.
\end{equation}
A questo punto le sezioni d'urto nucleari multigruppo per ciascun FSR sono 
richieste come input per il procedimento computazionale MOC e gli integrali 
mediati sull'area sono calcolati attraverso un pre processamento effettuato 
solitamente con metodi MonteCarlo.

\subsection{Soluzione mediante il calcolo del fattore di integrazione}
Ogni caratteristica può essere discretizzata in segmenti su FRSs individuali. 
Questa approssimazione permette all'equazione di poter essere localizzata sul 
segmento caratteristico $k$ lungo la FRS i-esima dal punto in cui entra $s'$ al 
punto in cui esce $s''$. Tale procedimento permette la valutazione analitica 
dell'integrale  attraverso la definizione di un fattore integrativo in termini 
di lunghezza ottica
\begin{equation}
\tau_{k,i,g}=\Sigma_{i,g}^T(s''-s') \,,
\end{equation}
ed esprimere il flusso uscente lungo la caratteristica come
\begin{equation}
\psi_{k,g}(s'')=\psi_{k,g}(s')e^{-\tau_{k,i,g}}+\frac{Q_{i,g}}{\Sigma_{i,g}^T}
(1-e^{-\tau_{k,i,g}}) \,.
\end{equation}
Con qualche passaggio è possibile ricavare la variazione del flusso angolare lungo la 
caratteristica 
\begin{equation}
\Delta\psi_{k,g}=\psi_{k,g}(s')-\psi_{k,g}(s'')=\biggl(\psi_{k,g}(s')-\frac{Q_{i
,g}}{\Sigma_{i,g}^T}\biggr)(1-e^{-\tau_{k,i,g}}) \,.
\end{equation}
La quantità chiave che rimane da determinare è l'integrale sopra l'area FSRs del 
flusso scalare $\phi_{i,g}(s)$ in \eqref{flusso_scalare_FSR}. Prima di tutto si 
definisce $l_{k,i}=s''-s'$ in modo che il flusso angolare medio nella FSR 
i-esima lungo la caratteristica $k$ sia espresso secondo il seguente integrale
\begin{equation}
\bar{\psi}_{k,i,g}=\frac{1}{l_{k,i}}\int_{s'}^{s''}ds\, \psi_{k,i,g}(s) \,.
\end{equation}
Per la valutazione dell'integrale, il flusso angolare medio lungo la 
caratteristica è espresso secondo la seguente espressione algebrica
\begin{equation}
\bar{\psi}_{k,i,g}=\frac{1}{l_{k,i}}\biggl[\frac{\psi_{k,g}(s')}{\Sigma_{i,g}^T}
(1-e^{-\tau_{k,i,g}})+\frac{l_{k,i}\,Q_{i,g}}{\Sigma_{i,g}^T}\biggl(1-\frac{
(1-e^{-\tau_{k,i,g}})}{\tau_{k,i,g}}\biggr)\biggr] \,.
\label{flusso_angolare_medio}
\end{equation}
Nel caso in cui sia le sezioni d'urto sia la sorgente siano costanti nella FSR 
i-esima, il valore del flusso angolare medio $\bar{\psi}_{k,i,g}$ data dalla 
\eqref{flusso_angolare_medio} risulta esatto.

\subsection{Il metodo di probabilità di collisione \label{probcoll}}
Il metodo di probabilità di collisione (CP) si ricava da una discretizzazione 
spaziale dell'equazione integrale in forma multigruppo, assumendo la sorgente di 
neutroni isotropa.

Per un problema contenente $I$ regioni, questo approccio produce una matrice $I 
\times I$ per ogni gruppo di energia. 
Questo metodo è preferito per trattare mesh generali non strutturate e problemi 
con poche regioni.
Le probabilità di collisione possono essere definite sopra un dominio infinito 
(come un reticolo di celle identiche o di assembly) o sopra un dominio finito (o 
aperto) $D$ contornato dalla superficie $\partial D$.
Il formalismo nel primo caso è semplice ma porta ad alcune difficoltà relative alla 
valutazione pratica del CPs. Nel secondo caso invece è richiesta l'aggiunta 
delle condizioni al contorno per chiudere il dominio.
Prima di tutto si integra sopra l'angolo solido l'equazione \eqref{soluz_caratteristica} in forma multigruppo per ottenere direttamente il flusso 
$\phi_g(\vec{r})$
\begin{equation}
\phi_g(\vec{r})=\int_{4\pi} d^2\Omega \, 
\phi_g(\vec{r},\vec{\Omega})=\frac{1}{4\pi}\int_{4\pi} d^2\Omega \, 
\int_0^\infty ds \, e^{-\tau_g(s)} Q_g(\vec{r}-s\vec{\Omega}) \,,
\end{equation}
dove $\tau_g$ è la lunghezza ottica in forma multigruppo definita in Eq. 
\eqref{cammino_ottico}.
Si introduce un cambio di variabile $\vec{r}'=\vec{r}-s\vec{\Omega}$ con 
$d^3r'=s^2\,d^2\Omega ds$ e si ottiene 
\begin{equation}
\phi_g(\vec{r})=\frac{1}{4\pi}\int_\infty d^3r' \, 
\frac{e^{-\tau_g(s)}}{s^2}Q_g(\vec{r}') \,,
\end{equation}
con $s=|\vec{r}-\vec{r}'|$.

Questa forma dell'equazione del trasporto è generalmente usata per rappresentare 
un reticolo infinito di celle identiche, che si ripetono e vengono modellate
mediante  condizioni al contorno periodiche. Successivamente si ottimizza la 
partizione delle celle unitarie o assembly all'interno delle regioni $V_i$.
Si rappresenta con il simbolo $V_i^\infty$ un infinito insieme di di regioni 
$V_i$ appartenenti a tutte le celle o all'assembly del reticolo. Inoltre si 
suppone che la sorgente secondaria di neutroni sia uniforme e uguale a $Q_{i,g}$ 
in ogni regione $V_i$. 
Dopo la moltiplicazione di $\Sigma_g(\vec{r})$ e integrato sopra ogni regione 
$V_i$, si ottiene
\begin{equation}
\int_{V_j}d^3r \, \Sigma_g(\vec{r})\phi_g(\vec{r})=\frac{1}{4\pi}\int_{V_j}d^3r 
\, \Sigma_g(\vec{r})\sum	_i Q_{i,g}\int_{V_i^\infty} d^3r' 
\frac{e^{-\tau_g(s)}}{s^2} \,.
\end{equation}
Tale equazione è semplificabile come
\begin{equation}
\label{flussocps}
V_j\Sigma_{j,g}\phi_{j,g}=\sum_i Q_{i,g}V_iP_{ij,g} \,,
\end{equation}
dove il pedice $j$ corrisponde all'indice dell'isotopo fissionabile, mentre
\begin{equation}
\phi_{j,g}=\frac{1}{V_j}\int_{V_j} d^3r\,\phi_g(\vec{r}) \,,
\end{equation}
\begin{equation}
\Sigma_{j,g}=\frac{1}{V_j\phi_{j,g}}\int_{V_J}d^r\,\Sigma_g(\vec{r})\phi_g(\vec{
r}),
\end{equation}
\begin{equation}
P_{ij,g}=\frac{1}{4\pi 
V_i}\int_{V_i^\infty}d^3r'\int_{V_j}d^3r\,\Sigma_g(\vec{r})\frac{e^{-\tau_g(s)}}
{s^2} \,.
\end{equation}

La collisione di probabilità $P_{ij,g}$ è la probabilità che un neutrone nato 
uniformemente e isotropicamente in una delle regioni $V_i$ del reticolo subisca 
la sua prima collisione nella regione $V_j$ della cella o assembly.

Se la sezione d'urto totale $\Sigma_g(\vec{r})$ è costate e uguale a 
$\Sigma{j,g}$ nella regione $V_j$, si può definire la probabilità di collisione 
ridotta come segue
\begin{equation}
p_{ij,g}=\frac{P_{ij,g}}{\Sigma{j,g}}=\frac{1}{4\pi 
V_i}\int_{V_i^\infty}d^3r'\int_{V_j}d^3r\,\frac{e^{-\tau_g(s)}}{s^2} \,.
\end{equation}
Le CPs ridotte generalmente rimangono finite nel limite di $\Sigma_{j,g}$ che
tende a zero. Questo permette un corretto comportamento della teoria di 
probabilità di collisione nel caso in cui regioni del reticolo siano considerate 
vuote.

Si possono scrivere inoltre le proprietà di reciprocità e di conservazione delle 
CPs, rispettivamente
\begin{equation}
p_{ij,g}V_j=p_{ji,g}V_j \,,
\end{equation}
\begin{equation}
\sum_j p_{ij,g}\Sigma_{j,g}=1; \quad \forall i \,.
\end{equation}
Usando la proprietà di reciprocità, l'equazione \eqref{flussocps} si semplifica 
ulteriormente
\begin{equation}
\label{flussocpssemplificato}
\phi_{i,g}=\sum_j Q_{j,g}p_{ij,g} \,.
\end{equation}
Il metodo di probabilità di collisione procede generalmente per tre passaggi.
Il primo è l'applicazione di un processo di tracking sopra la geometria del 
reticolo in modo da calibrare un numero sufficientemente grande di traiettorie 
dei neutroni. Successivamente è richiesto un processo di integrazione numerica, 
usando le informazioni di tracking e la conoscenza delle sezione d'urto 
macroscopiche totali per ogni regione. Questa integrazione deve essere fatta per 
ogni gruppo di energia, ma il processo è completamente parallelizzabile, in 
quanto non ci sono interazioni tra i gruppi. Una volta conosciuti i CPs, il 
flusso integrato può essere calcolato dalle Equazioni \eqref{flussocps} e 
\eqref{flussocpssemplificato}.
Infine le tecniche di probabilità di collisione possono anche essere applicate ad 
un caso con dominio $D$ contornato dalla superficie $\partial D$. 

